% (comment) First text note with sample.cls (this document class controls document formatting)
\documentclass{sample}
\begin{document}
This is where the document begins. You may need to also view this example as a source (tex) file.

LaTeX ignores blank/white spaces and justifies the document itself, placing \textit{interword} and \textit{intersentence} spaces automatically. 

To force an interword space, use a lone backslash \textbackslash \textbackslash, so the following \ \ \ \ is four spaces away. 

One can also introduce whitespace with the tilde \~\ character (this will attempt to keep the words on the same line). These~~words~~are~~forced~~to~~be~~separated~~by~~an~~extra~~space.

In LaTeX, one can document (escape) special characters with backslash. For example, comments start with \% (use \textbackslash textbackslash \% to document the back slash and \%).

LaTeX utilises standard commands, some with parameters. To document to command use the escape sequence \textbackslash command. To invoke a command, e.g. start a new paragraph, use \textbackslash par. (Alternatively, start a new paragraph with two carriage returns.)

Some common comands with parameters: \textbackslash textit\{\} for \textit{italicising}, \textbackslash textbf\{\} for \textbf{boldface} and \textbackslash texttt\{\} for \texttt{typewriter style} text.

Single left and right quotes are handled by ` and ' respectively, `for example'. Double quotes are handled by two left-single ` quotes  and then two right-single quotes (apostrophes) ', ``for example''.
\end{document}